\section{Literature Review}
\label{sec:literature}
\subsection{Structural IVQR and identification}
Quantile regression provides a way to study relationships across a conditional outcome distribution rather than compressing them into a conditional mean \parencite{koenker_bassett_1978}. Giving such relationships a structural treatment interpretation requires additional assumptions. When treatment is endogenous, observations with the same controls but different treatment values can differ in unobserved determinants of the outcome. A conditional quantile coefficient can consequently combine a treatment response with selection. Structural IVQR addresses this problem through excluded instruments and restrictions on structural ranks, rather than by treating an observational conditional quantile as a potential-outcome quantile \parencite{chernozhukov_hansen_2005,chernozhukov_hansen_2013}.

Rank invariance and rank similarity are distinct restrictions. The former maintains an individual's structural rank across treatment states, whereas the latter allows ranks to change subject to the relevant conditional distributional restriction. Their role is to connect the observed outcome to structural quantile moments. They do not by themselves ensure that these moments separate every competing treatment coefficient. The availability of valid instruments and the uniqueness or precision of identification are therefore separate questions \parencite{chernozhukov_hansen_2005,chernozhukov_hansen_wuthrich_2017}. This distinction matters when a computational criterion is evaluable over many candidate coefficients but changes little over the investigated domain.

\textcite{chernozhukov_hansen_2006} develop estimation and inference for the structural quantile process. Profiling nuisance coefficients at candidate treatment effects reduces the dimension of the difficult search, but it does not make the identifying information independent of the candidate or quantile. The present experiment retains that candidate-specific perspective. Its ordinary least-squares relevance diagnostics order the designed treatment--instrument relationship; they are not substitutes for the nonlinear structural identification conditions. The study thus asks about performance along a specified relevance path, not whether a linear first-stage threshold certifies IVQR identification.

\subsection{Weak instruments and IVQR inference}
The weak-instrument literature distinguishes uncertainty about a structural effect from the reliability of a particular approximation to its sampling distribution. \textcite{stock_wright_2000} show how weak identification changes the behavior of GMM estimation and testing; \textcite{andrews_stock_sun_2019} review the implications for linear-IV practice. These results motivate attention to candidate-specific tests, but their diagnostic thresholds and guarantees cannot simply be assigned to a nonlinear quantile procedure. The relevant moment conditions, nuisance estimates and reference distribution must be specified for the procedure actually used.

In IVQR, \textcite{chernozhukov_hansen_2008} develop robust candidate-based inference, and \textcite{jun_2008} study identification-robust tests and power. Robustness to identification quality is a property of an inferential procedure under its maintained conditions, not a promise that false candidates will be rejected frequently. A set can retain the true parameter with the intended probability and still be too broad to distinguish substantively different effects. Conversely, a sharp-looking profile is not evidence that its test is correctly calibrated. These established distinctions motivate joint reporting of retention and discrimination rather than a claim that weak-instrument power loss is a new finding.

Related approaches also differ in the distributional justification they require. \textcite{chernozhukov_hansen_jansson_2009} exploit a conditional pivotal property for finite-sample inference. \textcite{andrews_mikusheva_2016} develop conditional inference with a functional nuisance parameter, allowing information in the candidate-moment process to enter testing. Neither approach is implemented by merely comparing a historical statistic with a chi-square quantile. Our comparisons therefore describe the retained reference-cutoff rule. They do not implement these alternatives, establish equal-size power rankings or infer a uniform guarantee from nominal-looking acceptance frequencies.

\subsection{High-dimensional quantile regression and orthogonalization}
Many controls create a nuisance-estimation problem even when the structural parameter of interest is low dimensional. Approximate sparsity provides one way to make that problem manageable: a limited collection of terms can approximate the relevant nuisance function despite a much larger set of candidates. \textcite{belloni_chernozhukov_2011} develop penalized quantile regression and a pivotal penalty choice under suitable conditions. Such regularization introduces its own approximation and estimation errors; successful prediction or variable selection is not itself an inference argument for a treatment coefficient.

Orthogonal scores address the propagation of nuisance errors into the target moment. \textcite{belloni_chernozhukov_kato_2019} develop post-selection inference for approximately sparse quantile models, while \textcite{chernozhukov_et_al_2018} provide a general double/debiased machine-learning framework. Orthogonality is a local derivative property of a population moment. Cross-fitting is a sampling arrangement that separates nuisance fitting from score evaluation. They serve related but different purposes, and neither turns an uninformative instrument into an informative one. Moreover, an orthogonality calculation requires its own differentiability and moment conditions; it cannot be invoked solely because the executable procedure uses a regularized regression.

These distinctions are important here because the retained procedures use same-sample nuisance fits, a fitted-normal density proxy and different projection conventions. The experiment compares the complete implementations, rather than changing only the presence of a penalty. The ideal construction in Section~\ref{sec:framework} explains the motivation; Appendix~\ref{app:reproducibility} identifies the actual conventions. The non-finite mixed residual density limits the available standard justification without making the observed finite-sample summaries undefined. Absence of cross-fitting alone is likewise insufficient to diagnose the cause of an observed rejection pattern.

\subsection{DML-IVQR simulations, computation and the remaining gap}
\textcite{chen_huang_tien_2021} supply the direct high-dimensional simulation comparator: they compare point errors across Oracle-GMM, Full-GMM and DML-IVQR, illustrate candidate-statistic profiles, and examine tuning and cross-fitting choices. Their displayed treatment assignment holds excluded-signal coefficients fixed. Our contribution is a controlled extension of that experiment, not the introduction of profiling or regularized IVQR. Preserving the executable lineage also requires acknowledging that paper formulas, source-code conventions and every published numerical comparison are not interchangeable objects.

Computational IVQR research addresses additional questions. \textcite{chen_lee_2018} formulate exact GMM computation using mixed-integer optimization; \textcite{kaido_wuthrich_2021} decompose estimation into more tractable quantile-regression subproblems. \textcite{kaplan_sun_2017} study smoothed estimating equations, including their computational and inferential consequences. These approaches show why the numerical route deserves explicit attention, but none is a substitute calculation performed in the present archive. Changing the objective, solver or reference rule would define a different comparison and could not silently replace the frozen historical results.

The neighboring contribution of \textcite{jin_sun_2025} concerns an orthogonal moment for an unconditional quantile-treatment-effect target with high-dimensional covariates. It should not be conflated with an implemented comparison of the present structural coefficient and candidate grid. More generally, differences in target, nuisance dimension and algorithm restrict what can be inferred from comparisons across papers. The gap claimed here is consequently specific to the verified predecessor experiment, rather than a claim that no earlier work varied identification quality or combined several performance measures.

The current experiment follows accuracy, direct truth acceptance, specified false-candidate rejection and finite-grid localization along a relevance modification that preserves the documented joint treatment/outcome-disturbance distribution. This organization makes disagreement among the diagnostics visible. It is consistent with evaluating simulation methods against explicitly stated aims and performance measures \parencite{morris_white_crowther_2019}. The resulting contribution is experimental and diagnostic: it documents where apparently reassuring summaries fail to describe discrimination in this DGP family. It supplies neither a new estimator nor an inference theorem for the preserved procedures.
